\section{Natural and upper densities separate in every number field}
\label{sec:ideal-separation}

This section extends the smooth-factor optimization in the incoming
\path{grh_transfers_density_gaps_and_zero_tests.zip}~\cite{prior-density-gaps} from integers
and polynomials to ideals of an arbitrary number field.  Its additional
ingredient is a single prime-power obstruction that works in every field,
including when distinct ideals have the same norm.  All assertions in this
section are unconditional.

Let $K$ be a number field of degree $d$, and let $\mathcal I_K$ be its
nonzero integral ideals.  Put
\[
 A_K(X)=\#\{\mathfrak a\in\mathcal I_K:N\mathfrak a\le X\}.
\]
We use the classical ideal-counting theorem
$A_K(X)=\kappa_K X+o(X)$, where $\kappa_K>0$; see
\cite{murty-vanorder}.  Natural density means a limit of the
proportion in this norm ball, and upper density means its limsup.

Fix either of two conventions.  In the \emph{integral} convention the
forbidden triples are
$\mathfrak a,\mathfrak a\mathfrak r,
\mathfrak a\mathfrak r^2$ with integral
$\mathfrak r\ne\mathcal O_K$.  In the \emph{fractional} convention
$\mathfrak r\ne\mathcal O_K$ may be any nonzero fractional ideal, as long
as all three terms are integral.  The latter includes nontrivial ratios
of norm one.  In valuation coordinates these are, respectively,
nonconstant arithmetic progressions with a nonnegative step and with an
arbitrary integral step.  Write $b_K$ for the supremum of natural
densities of admissible sets, and $u_K$ for the supremum of upper
densities, using the same fixed convention in both definitions.

\begin{theorem}[Separation and finite approximation]
\label{thm:ideal-separation}
For every number field and for either convention,
\[
                         u_K>b_K.
\]
Both numbers are computable from an explicit presentation of $K$.
More precisely, fix a prime ideal $\mathfrak p$, put
$Q=N\mathfrak p$ and $M=Q^8$.  For every integer $Y\ge Q$, define
\[
 c_Y=\prod_{N\mathfrak q\le Y}(1-(N\mathfrak q)^{-1}).
\]
Then
\begin{equation}\label{eq:ideal-gap}
 u_K-b_K\ \ge\ \frac{c_Y}{M(M+1)}-\frac{d}{Y}.
\end{equation}
The right side is positive for every sufficiently large $Y$.
\end{theorem}

\subsection{The two local optimization problems}

Let $S$ be the prime ideals of norm at most $Y$.  An ideal is
$S$-smooth if it has no prime factor outside $S$, and $S$-rough if it
has none in $S$.  Every ideal has a unique factorization
$\mathfrak a=\mathfrak g\mathfrak h$ with rough $\mathfrak g$ and
smooth $\mathfrak h$.  Set
\[
 c=\prod_{\mathfrak q\in S}(1-(N\mathfrak q)^{-1}),\qquad
 v=\prod_{\mathfrak q\notin S}(1-(N\mathfrak q)^{-2}).
\]
For admissible subsets $D$ of the smooth ideals, and for $t\ge1$, put
\begin{align}
 B&=c\sup_D\sum_{\mathfrak h\in D}\frac1{N\mathfrak h},\\
 J(t)&=\max\{\#D:D\text{ admissible},\ N\mathfrak h\le t
                         \text{ for all }\mathfrak h\in D\},\\
 U&=c\int_1^\infty J(t)\,\frac{dt}{t^2}.
\end{align}
The smooth reciprocal sum is exactly $1/c$.  In particular all these
expressions converge, $B\le U\le1$, and the weighted supremum defining
$B$ is attained.  For the latter assertion one may use compactness of
the space of subsets of the countable smooth monoid: admissibility is
closed and the summable weighted objective is continuous.

\begin{lemma}[Local optimization bounds]\label{lem:BU}
For either ratio convention,
\[
                  vB\le b_K\le B,\qquad vU\le u_K\le U.
\]
\end{lemma}
\begin{proof}
Finite inclusion--exclusion and ideal counting give rough ideals
relative density $c$.  Squarefree rough ideals have relative density
$cv$.  Indeed, impose squarefreeness at finitely many further primes
and then let that finite set increase.  The omitted-square error is
bounded by a constant times $\sum(N\mathfrak q)^{-2}$, uniformly as
$X\to\infty$, since $A_K(t)\ll_K t$.  This convergent bound justifies
passage to the infinite product.

For the upper bound on $b_K$, let $E$ have natural density $b$.  Partial
summation yields
\[
 \sum_{\substack{\mathfrak a\in E\\N\mathfrak a\le X}}
       \frac1{N\mathfrak a}
   =(b\kappa_K+o(1))\log X.
\]
Each rough fiber $D_{\mathfrak g}$ is admissible.  Consequently its
smooth reciprocal sum is at most $B/c$, and the full sum is bounded by
\[
 \frac Bc\sum_{\substack{\mathfrak g\text{ rough}\\N\mathfrak g\le X}}
          \frac1{N\mathfrak g}
   =(B\kappa_K+o(1))\log X.
\]
Thus $b\le B$.

For the upper bound on $u_K$, count in each rough fiber at cutoff $X$:
\[
 \#\{\mathfrak a\in E:N\mathfrak a\le X\}
 \le \sum_{\substack{\mathfrak g\text{ rough}\\N\mathfrak g\le X}}
                  J(X/N\mathfrak g).
\]
Rough ideal counting, first on a bounded interval $1\le t\le T$, gives
limit $c\int_1^T J(t)t^{-2}\,dt$ after division by $A_K(X)$.
The tail is uniformly controlled by the corresponding smooth counting
function $H(t)$: $H(t)=O_S((1+\log t)^{|S|})$ and
$\int_1^\infty H(t)t^{-2}\,dt=1/c$.  Equivalently, expand by smooth
ideals and dominate each term by $O_K(1/N\mathfrak h)$.
Letting $T\to\infty$ proves the upper bound.

For both lower bounds retain only squarefree rough parts.  In any
progression, each valuation outside $S$ would form an arithmetic
progression in $\{0,1\}$, so it is constant.  All three terms must
therefore lie in a single rough fiber.  Using the same admissible
smooth set $D$ in every such fiber gives natural density
$cv\sum_{\mathfrak h\in D}1/N\mathfrak h$.  Taking the supremum
proves $b_K\ge vB$.

To attain the lower upper-density bound, choose $X_j$ with
$X_j/X_{j-1}\to\infty$.  For every squarefree rough $\mathfrak g$
with $X_{j-1}<N\mathfrak g\le X_j$, permanently assign a maximizing
smooth set through $X_j/N\mathfrak g$.  This is one infinite admissible
set, since each fiber is assigned once.  At cutoff $X_j$ and fixed
$T$, its newly assigned fibers with $X_j/T<N\mathfrak g\le X_j$
give limiting proportion $cv\int_1^T J(t)t^{-2}\,dt$.
Let $T\to\infty$ to obtain upper density at least $vU$.
\end{proof}

\subsection{A uniform obstruction to compatible prefix optima}

Let $A$ be the nonnegative integers with ternary digits only $0,1$.
The greedy smooth set $G$ consists of ideals all of whose prime
exponents lie in $A$.  The following proof also resolves the issue of
equal norms: a norm layer need not be either entirely included or
entirely excluded.

For $n=\sum_j d_j3^j$, define
\[
 a(n)=\sum_{d_j=1}3^j,\qquad
 b(n)=\sum_{d_j\in\{1,2\}}3^j.
\]
Then $a(n),b(n)\in A$ and $n=2b(n)-a(n)$.  If $n\in A$ both equal
$n$; otherwise $0\le a(n)<b(n)<n$.  Applying this coordinatewise
shows that every nongreedy ideal is the final term of a progression
whose first two terms are greedy ideals of strictly smaller norm.
The ratio is integral.  Meanwhile all greedy ideals coexist without
any progression, even with fractional ratios.

It follows by induction on integer norm that any admissible set
simultaneously maximizing the number of ideals at \emph{every} cutoff
$1,2,\ldots,M$ must equal $G$ through $M$.  At a new norm layer every
nongreedy candidate is already blocked, and all greedy candidates can
be included together.  To achieve even the greedy count one must
include every such candidate.  This argument does not choose an
ordering within a norm layer.

Now change only the pure powers of the selected prime $\mathfrak p$
at cutoff $M=Q^8$.  Replace exponent set
\[
                    \{0,1,3,4\}\quad\text{by}\quad
                    \{0,1,3,7,8\}.
\]
The latter is arithmetic-progression-free and has one additional
element.  The resulting smooth set is admissible even for fractional
ratios.  To check this, any new progression has an added pure
$\mathfrak p$-power as one term.  At every other prime one of its
three valuations is zero; the other two belong to $A$ (or are also
zero).  Since $0\in A$ and $A$ is progression-free, all three are
zero.  The progression is confined to pure $\mathfrak p$-powers,
where the new exponent set forbids it.

Thus $G$ is not cardinality-optimal through $M$.  No admissible set
can optimize all prefixes through $M$.  For every admissible smooth
$D$ there is some integer $1\le k\le M$ with
$J(k)-\#\{\mathfrak h\in D:N\mathfrak h\le k\}\ge1$.
That deficit persists throughout $[k,k+1)$, giving
\[
 \int_1^\infty
 \bigl[J(t)-\#\{\mathfrak h\in D:N\mathfrak h\le t\}\bigr]
          \frac{dt}{t^2}
 \ge\frac1{k(k+1)}\ge\frac1{M(M+1)}.
\]
Taking the infimum over $D$ proves
\begin{equation}\label{eq:BUgap}
                         U-B\ge\frac{c}{M(M+1)}.
\end{equation}

\subsection{Passing to the global densities}

There are at most $d$ prime ideals of any specified integer norm:
that norm must be a power of a unique rational prime, and at most $d$
primes lie over it.  Therefore
\[
 1-v\le\sum_{N\mathfrak q>Y}(N\mathfrak q)^{-2}
       \le d\sum_{n>Y}n^{-2}\le d/Y.
\]
Equations~\eqref{eq:BUgap} and Lemma~\ref{lem:BU} now give
\[
 u_K-b_K\ge vU-B\ge U-B-(1-v)
                     \ge c_Y/[M(M+1)]-d/Y.
\]
Finally
\[
 c_Y\ge\prod_{p\le Y}(1-p^{-1})^d\gg_d(\log Y)^{-d}
\]
by the classical Mertens product theorem.  This decreases more slowly
than $d/Y$, proving strict positivity.  One could instead use the
number-field Mertens theorem, but it is unnecessary here.

For computability choose $Y$ with $d/Y$ below the desired tolerance.
The finite smooth reciprocal sum through $H$ approaches $1/c$, with
explicit rational tail
$1/c-\sum_{N\mathfrak h\le H}1/N\mathfrak h$.
Finite weighted independent-set optimization through $H$ approximates
$B$ with at most $c$ times this tail.  Optimizing cardinalities at
each norm through $H$ and integrating the resulting step function
approximates $U$; the tail is bounded by
$c\int_H^\infty H(t)t^{-2}\,dt$.  Writing $H_S(t)$ for the smooth
counting function to avoid confusion with the cutoff, this integral is
the computable rational number
\[
 c\left(\frac{H_S(H)}{H}+\frac1c
             -\sum_{N\mathfrak h\le H}\frac1{N\mathfrak h}\right).
\]
It tends to zero as $H$ increases. Lemma~\ref{lem:BU} then
computes $b_K$ and $u_K$.  This is an existence-of-algorithm result;
the exact independent-set problems can be expensive.

\subsection{An explicit certificate in the conductor-23 field}

For $K=\mathbb Q(\sqrt{-23})$, choose a prime of norm $2$, so $M=256$.
The script \path{code/check_ideal_density.py} computes $c_Y$ at
$Y=10^7$ using only integer arithmetic.  For each rational prime,
the splitting test is its quadratic-residue class modulo $23$;
$23$ ramifies, split primes contribute two factors, and inert primes
contribute one factor when $p^2\le Y$.

After each Euler factor the script rounds \emph{down} on the grid
$10^{-40}\mathbb Z$.  Its certified lower bound is
\[
 c_{10^7}\ge
 \frac{11078476589752147709803027083658225201}
      {625000000000000000000000000000000000000}.
\]
There are $664449$ factors.  Using the exact tail bound $2/10^7$
in~\eqref{eq:ideal-gap} proves, under either convention,
\begin{equation}\label{eq:ideal-explicit-gap}
                         u_K-b_K>\frac1{20000000}.
\end{equation}
The sharper rational lower bound is stored in the JSON certificate.
This computation checks a finite numerical part of the theorem; the
norm-counting limits and infinite construction are established by the
proof above.

For fractional ratios, the theorem implies upper density strictly
larger than greedy density, since the greedy set has a natural
density at most $b_K$.  It does \emph{not} show that the fractional-ratio
natural-density supremum exceeds greedy density.  Thus it contributes
to the extensions requested in Questions 5.3 and 5.5 of
Best et al.~\cite{best}, while their Question 5.4 remains unresolved
by this report.  No assertion of priority for the general separation
theorem is made.
